% Визуальные стили книги: цвета, типы табличных колонок,
% базовый стиль notebox и команды нотификационных блоков.
% Преамбулу не трогать при изменении оформления — править здесь.

\usepackage[table]{xcolor}
\usepackage[skins,breakable]{tcolorbox}
\usepackage{listings}
\usepackage{listingsutf8}
\tcbuselibrary{listings}

% --- Цвета фигур (используются в Inkscape-схемах и таблицах) ---
\definecolor{Ink}{HTML}{1A2030}
\definecolor{Soft}{HTML}{F5F7FD}
\definecolor{Panel}{HTML}{E8ECF7}
\definecolor{AccentA}{HTML}{4E63D9}
\definecolor{AccentB}{HTML}{7759D6}
\definecolor{Muted}{HTML}{5C647A}
\definecolor{Good}{HTML}{2F8B67}
\definecolor{Warn}{HTML}{B8821C}
\definecolor{Bad}{HTML}{C15462}

% --- Цвета notebox-блоков и текстового маркера ---
\definecolor{ImportantInk}{HTML}{A1273C}
\definecolor{PracticeInk}{HTML}{1C56A3}
\definecolor{MarkerYellow}{HTML}{FFF2A8}

% --- Пастельная палитра купюр Банка России (верхний градиент на старте части) ---
% Соответствие частей сохранено: I→100, II→200, III→1000, IV→2000, V→5000.
% Тона смягчены до пастели: прежние насыщенные купюрные цвета давали жёсткую
% полосу — здесь их светлые версии для диффузного mesh-градиента (ниже).
\definecolor{BnHundred}{HTML}{EAD9B4}      % 100  — песочный  (был горчично-золотой)
\definecolor{BnTwoHundred}{HTML}{C6E6CC}   % 200  — мятный    (был салатовый)
\definecolor{BnThousand}{HTML}{BFE5DC}     % 1000 — аква       (был бирюзовый)
\definecolor{BnTwoThousand}{HTML}{BFD7EC}  % 2000 — небесный   (был синий)
\definecolor{BnFiveThousand}{HTML}{EDC8BC} % 5000 — пудровый   (был коралл)

% --- Мягкий диффузный градиент на странице с заголовком части ---
% Вместо вертикальной полосы «цвет→белый» — два пастельных пятна (bloom),
% растворяющихся в белом поле страницы (path fading даёт прозрачность к краю):
% насыщенное вверху-справа за заголовком \part, светлее — слева и ниже; центр
% и верх-слева остаются почти белыми. Жёсткой границы полосы нет.
% Вызов \introgradient{<цвет>} ставится первой строкой в src/parts/partN/intro.tex —
% AddToShipoutPictureBG* применяет фон к странице, на которой выводится \part-
% заголовок и начало intro (после переопределения \@endpart разрыва нет).
\usepackage{tikz}
\usetikzlibrary{fadings}
\tikzfading[name=introbloom, inner color=transparent!6, outer color=transparent!100]
\newcommand{\introgradient}[1]{%
  \AddToShipoutPictureBG*{%
    \begin{tikzpicture}[remember picture, overlay]
      \coordinate (NW) at (current page.north west);
      \coordinate (NE) at (current page.north east);
      % основное пятно — вверху, смещено вправо (ложится за заголовок \part)
      \fill[#1, path fading=introbloom]
      ([xshift=-16mm, yshift=-28mm] NE) circle (96mm);
      % второе пятно — слева, светлее и чуть ниже (диагональ mesh)
      \fill[#1!58, path fading=introbloom]
      ([xshift=20mm, yshift=-72mm] NW) circle (78mm);
    \end{tikzpicture}%
  }%
}

% --- Запрет переноса по слогам внутри таблиц ---
% Применяется через >{...}-префикс в Y, P, B-колонках. Авто-перенос
% (по словарю слогов) выключен; перенос по уже стоящим в тексте дефисам
% и слешам остаётся разрешённым, иначе длинные идентификаторы вроде
% `ReadTransactionsBasic` дают огромные overfull. \emergencystretch
% из preamble.tex продолжает спасать в~остальных случаях.
\newcommand{\tblnohyphens}{\hyphenpenalty=10000\relax}

% --- Колонки для tabularx ---
% Y — растягивающаяся ragged-right без переносов
% P{w} — фиксированная ширина, ragged-right, без переносов
% B{w} — то же + \bfseries (для левого «label»-столбца сравнительных таблиц)
\newcolumntype{Y}{>{\tblnohyphens\RaggedRight\arraybackslash}X}
\newcolumntype{P}[1]{>{\tblnohyphens\RaggedRight\arraybackslash}p{#1}}
\newcolumntype{B}[1]{>{\tblnohyphens\RaggedRight\arraybackslash\bfseries}p{#1}}

% --- Стиль таблиц (Packt header-strip) ---
% Тёмная плашка Ink на всю ширину строки хедера, белый PT~Sans~Bold в~ячейках.
% Тело таблицы — белое, без зебры. Booktabs-линейки.
%
% Использование:
%   \begin{tabularx}{\textwidth}{@{} B{0.18\textwidth} Y Y Y Y @{}}
%     \toprule
%     \headrow \thd{Параметр} & \thd{Visa} & \thd{Mastercard} & ... \\
%     \midrule
%     Базовый профиль & ... \\
%     ...
%     \bottomrule
%   \end{tabularx}
\newcommand{\headrow}{\rowcolor{Ink}}
\newcommand{\thd}[1]{{\color{white}\sffamily\bfseries #1}}

% --- Легенда-подпись под таблицей ---
% Нейтральная плашка для пояснения категорийных заливок (зелёный/красный/
% серый и т.~п.). Ставится между \end{tabularx} и \caption внутри table.
% Тон серый, чтобы не конкурировать с тёмным хедером и категорийными
% подсветками тела таблицы.
\newcommand{\tablenote}[1]{%
  \par\smallskip
  \fcolorbox{black!18}{black!3}{%
    \parbox{\dimexpr\linewidth-2\fboxsep-2\fboxrule\relax}{%
      \footnotesize\color{Muted}#1%
    }%
  }%
}

% --- Унификация табличного тела ---
% Канон набора тела любой таблицы: \footnotesize, tabcolsep 4pt,
% arraystretch 1.2, центрирование (нужно, когда ширина tabularx < \textwidth).
% Снимает копипасту \begingroup…\renewcommand, разбредавшуюся по 35 таблицам
% (tabcolsep гулял 2/3/4/6pt, stretch — 1.15–1.3), и даёт единую точку
% правки. Исключения (плотные широкие матрицы, разреженные пары) переопределяют
% \tabcolsep, \arraystretch или кегль отдельной строкой после \tablesetup,
% например: \tablesetup\setlength{\tabcolsep}{3pt}.
%
% ВАЖНО про центрирование. \centering центрирует tabularx только если абзац
% завершён \par'ом, пока группа ещё открыта; поэтому таблицу закрывать
% \end{tabularx}\par перед \endgroup. Без \par блок прижимается влево, вся
% пустота уходит вправо (для узкой таблицы это так же плохо, как растяжка
% на всю полосу). Обёртка-environment не годится: tabularx собирает тело,
% сканируя поток до литерального \end{tabularx}; спрятанный в end-коде
% окружения, он даёт «File ended while scanning». Канонический шаблон:
%   \begingroup\tablesetup
%   \begin{tabularx}{<доля>\textwidth}{<колонки>}
%     \toprule \headrow \thd{..} & .. \\ \midrule
%     строка \\ \rowsep строка \\ ... строка \\ \bottomrule
%   \end{tabularx}\par
%   \endgroup
\newcommand{\tablesetup}{%
  \footnotesize
  \setlength{\tabcolsep}{4pt}%
  \renewcommand{\arraystretch}{1.2}%
  \centering
}

% --- Волосяной разделитель строк таблиц ---
% Тонкая линейка black!16 между строками: горизонтальный ориентир, особенно
% для многострочных (переносящихся) построений, на которые жалуются читатели.
% Не конфликтует с booktabs-обвязкой (top/mid/bottom остаются чёрными) и с
% категорийными заливками ячеек — в отличие от зебры. Ставится после `\\`
% между ВСЕМИ строками тела каждой таблицы: единый принцип обязателен —
% выборочные линейки (часть таблиц с ними, часть без) читаются как
% непоследовательность, а не как осмысленный стиль. Отступы 0.3pt минимальны
% (на 20-строчной таблице +19pt высоты). Вес и тон правятся строкой ниже.
\newcommand{\rowsep}{\arrayrulecolor{black!16}\specialrule{0.4pt}{0.3pt}{0.3pt}\arrayrulecolor{black}}

% --- Текстовый маркер ---
\soulregister{\emph}{1}
\soulregister{\textit}{1}
\soulregister{\textbf}{1}
\soulregister{\texttt}{1}
\soulregister{\mbox}{1}
\soulregister{\enquote}{1}
% Кавычки внутри \highlight писать как \mbox{\enquote{...}}: soul красит
% содержимое \mbox целиком вместе с ёлочками, тогда как делитеры
% зарегистрированного \enquote без обёртки остаются вне жёлтого фона.
% \mbox запрещает перенос внутри закавыченного — для коротких терминов это норма.
\DeclareRobustCommand{\highlight}[1]{%
  \begingroup
  \sethlcolor{MarkerYellow}%
  \hl{#1}%
  \endgroup
}

% --- Базовый стиль notebox ---
% Хедер-плашка во всю ширину (белый sans-serif по цветному фону) + белое тело
% + тонкая цветная линейка снизу. Без рамки слева/справа/сверху — блок
% не «давит» на полосу и хорошо читается как самостоятельная единица.
\tcbset{
  notebox/.style={
    enhanced,
    breakable=false,
    boxrule=0pt,
    arc=0pt, outer arc=0pt,
    frame hidden,
    before skip=\medskipamount,
    after skip=\medskipamount,
    left=6pt, right=6pt, top=4pt, bottom=4pt,
    colback=white,
    fontupper=\rmfamily\footnotesize\linespread{0.95}\selectfont
    \setlength{\parindent}{0pt}%
    \setlength{\parskip}{0.5\baselineskip plus 1pt minus 1pt},
    parbox=false,
    fonttitle=\sffamily\footnotesize\bfseries,
    coltitle=white,
    titlerule=0pt,
    toptitle=3pt, bottomtitle=3pt,
    lefttitle=8pt, righttitle=8pt,
  },
}

% --- Цветной notebox с заголовком (общая основа) ---
% \colornotebox{<цвет>}{<заголовок>}{<содержимое>}
\newcommand{\colornotebox}[3]{%
  \begin{tcolorbox}[notebox,
      colbacktitle=#1,
      title={#2},
    borderline south={0.4pt}{0pt}{#1}]
    #3
\end{tcolorbox}}

% --- Тематические notebox-блоки ---
% Категория идёт капителью в плашке-хедере.
\newcommand{\importantnote}[1]{\colornotebox{ImportantInk}{ВАЖНО}{#1}}
\newcommand{\practicenote}[1]{\colornotebox{PracticeInk}{ПРАКТИКА}{#1}}

% --- Блок кода с подсветкой синтаксиса (listings + tcolorbox) ---
% Использование:
%   \begin{codeblock}{Python}
%   def f(x):
%       return x * 2
%   \end{codeblock}
% Внутрь аргумента notebox-команд код не поместить (verbatim не работает
% в аргументе макроса) — codeblock ставится рядом с поясняющей врезкой.
\definecolor{CodePanel}{HTML}{F4F6FB}
\definecolor{CodeKey}{HTML}{1C56A3}
\definecolor{CodeStr}{HTML}{2F8B67}
\definecolor{CodeCom}{HTML}{6B7280}
\definecolor{CodeNum}{HTML}{A1273C}

\lstdefinestyle{book-code}{
  basicstyle=\ttfamily\scriptsize\linespread{0.9}\selectfont,
  keywordstyle=\bfseries\color{CodeKey},
  commentstyle=\itshape\color{CodeCom},
  stringstyle=\color{CodeStr},
  numberstyle=\tiny\color{CodeCom},
  identifierstyle=\color{Ink},
  showstringspaces=false,
  tabsize=2,
  breaklines=true,
  breakatwhitespace=true,
  columns=fullflexible,
  keepspaces=true,
  upquote=true,
  extendedchars=true,
  inputencoding=utf8,
  literate={->}{{$\rightarrow$}}{2}
  {!=}{{$\neq$}}{2}
  {ë}{{\"e}}1{Ë}{{\"E}}1{š}{{\v{s}}}1{Š}{{\v{S}}}1
  {ž}{{\v{z}}}1{Ž}{{\v{Z}}}1{č}{{\v{c}}}1{Č}{{\v{C}}}1
  {ŝ}{{\^{s}}}1{Ŝ}{{\^{S}}}1{è}{{\`{e}}}1{È}{{\`{E}}}1
  {û}{{\^{u}}}1{Û}{{\^{U}}}1{â}{{\^{a}}}1{Â}{{\^{A}}}1
  {ʹ}{{$'$}}1{ʺ}{{$''$}}1
  {а}{{\cyra}}1{б}{{\cyrb}}1{в}{{\cyrv}}1{г}{{\cyrg}}1
  {д}{{\cyrd}}1{е}{{\cyre}}1{ё}{{\cyryo}}1{ж}{{\cyrzh}}1
  {з}{{\cyrz}}1{и}{{\cyri}}1{й}{{\cyrishrt}}1{к}{{\cyrk}}1
  {л}{{\cyrl}}1{м}{{\cyrm}}1{н}{{\cyrn}}1{о}{{\cyro}}1
  {п}{{\cyrp}}1{р}{{\cyrr}}1{с}{{\cyrs}}1{т}{{\cyrt}}1
  {у}{{\cyru}}1{ф}{{\cyrf}}1{х}{{\cyrh}}1{ц}{{\cyrc}}1
  {ч}{{\cyrch}}1{ш}{{\cyrsh}}1{щ}{{\cyrshch}}1{ъ}{{\cyrhrdsn}}1
  {ы}{{\cyrery}}1{ь}{{\cyrsftsn}}1{э}{{\cyrerev}}1{ю}{{\cyryu}}1
  {я}{{\cyrya}}1
  {А}{{\CYRA}}1{Б}{{\CYRB}}1{В}{{\CYRV}}1{Г}{{\CYRG}}1
  {Д}{{\CYRD}}1{Е}{{\CYRE}}1{Ё}{{\CYRYO}}1{Ж}{{\CYRZH}}1
  {З}{{\CYRZ}}1{И}{{\CYRI}}1{Й}{{\CYRISHRT}}1{К}{{\CYRK}}1
  {Л}{{\CYRL}}1{М}{{\CYRM}}1{Н}{{\CYRN}}1{О}{{\CYRO}}1
  {П}{{\CYRP}}1{Р}{{\CYRR}}1{С}{{\CYRS}}1{Т}{{\CYRT}}1
  {У}{{\CYRU}}1{Ф}{{\CYRF}}1{Х}{{\CYRH}}1{Ц}{{\CYRC}}1
  {Ч}{{\CYRCH}}1{Ш}{{\CYRSH}}1{Щ}{{\CYRSHCH}}1{Ъ}{{\CYRHRDSN}}1
  {Ы}{{\CYRERY}}1{Ь}{{\CYRSFTSN}}1{Э}{{\CYREREV}}1{Ю}{{\CYRYU}}1
  {Я}{{\CYRYA}}1,
}

\tcbset{
  codeblock-style/.style={
    enhanced, breakable,
    boxrule=0pt, arc=0pt, outer arc=0pt,
    frame hidden,
    before skip=\medskipamount, after skip=\medskipamount,
    left=10pt, right=6pt, top=4pt, bottom=4pt,
    colback=CodePanel,
    borderline west={2pt}{0pt}{Muted!50},
    listing only,
    overlay middle and last={%
      \node[anchor=north east,inner sep=2pt,font=\sffamily\tiny\color{CodeCom}]
      at ([xshift=-4pt,yshift=-2pt]frame.north east) {(продолжение)};
    },
  },
}

\newtcblisting{codeblock}[1]{%
  codeblock-style,
  listing options={style=book-code, language=#1},
}

% Подключить внешний файл с кодом, проверенный в samples/.
%   \codefile{Python}{samples/chXX-…/file.py}
\newcommand{\codefile}[2]{%
  \tcbinputlisting{%
    codeblock-style,
    listing engine=listings,
    listing options={style=book-code, language=#1},
    listing file={#2},
  }%
}

% Фрагмент внешнего файла с кодом по диапазону строк (см. samples/).
%   \codeexcerpt{Python}{samples/chXX-…/file.py}{firstline}{lastline}
\newcommand{\codeexcerpt}[4]{%
  \tcbinputlisting{%
    codeblock-style,
    listing engine=listings,
    listing options={style=book-code, language=#1, firstline=#3, lastline=#4},
    listing file={#2},
  }%
}
